\label{sec:intro_body}

Public question-and-answer platforms such as Stack Overflow are
part of the digital infrastructure of distributed
problem-solving: they convert individual software problems into
searchable, accepted, moderator-retained posts that future
contributors, users, and---increasingly---large language models
consult when facing analogous problems
\citep{wasko2005faraj,goes2014popularity,nagle2018learning}.
Generative AI reconfigures this infrastructure: since November
2022 a wide class of programming questions can be routed to a
private conversational LLM at near-zero marginal search cost
\citep{eloundou2024gpts,acemoglu2024simple}, and that private
channel produces no public artefact unless the user voluntarily
reposts the resolved problem. The substitution between the public
and private channels---and its consequences for a platform whose
replenishment depends on the contribution incentives the
substitute now competes against
\citep{kankanhalli2005contributing,butler2001membership,benkler2002coase}---is
the question this paper studies.

Recent work establishes that aggregate replenishment has slowed.
\citet{delriochanona2024large} report a $\sim$15\% mean decline
rising to $\sim$25\%; \citet{burtch2024chatgpt} document
analogous declines across Q\&A platforms; and
\citet{quinngutt2024heterogeneous} report that the effect of
generative AI on knowledge seeking is heterogeneous rather than
uniform, motivating a move beyond a single average effect to the
within-community margin. \citet{xue2026chatgpt}, in \emph{ISR},
use a cross-platform design and document both a volume reduction
and an \emph{improvement} in the average quality of surviving
questions---a combination naturally read as \emph{benign
pruning}, under which LLMs absorb the easier, lower-value
questions. A complementary ISR line studies \emph{answerer}
behaviour \citep{shanqiu2026isr}. The aggregate downward shift is
no longer in dispute; three questions remain open.

\emph{First}, the aggregate papers cannot adjudicate the
\emph{value distribution} of the displaced posts. The
benign-pruning reading requires that displacement concentrates at
the low-value end---that the platform retains the crowd-validated,
high-value answers downstream users rely on---which an aggregate
quantity-and-quality decomposition cannot test because its unit is
the average surviving question, not the score-conditioned post
count. \emph{Second}, none estimates the LLM-substitutability
mechanism at the \emph{task level}: substitutable categories
within AI-answerable tags should decline more than
non-substitutable categories \emph{in the same tags}, the slice
that admits a more disciplined quasi-causal reading under a
within-tag parallel-trends assumption; a cross-platform design
identifies the average treated-platform effect but cannot
decompose this share. \emph{Third}, existing studies measure
platform \emph{activity} but do not trace the same association
through \emph{funnel-qualified posts} (accepted, non-closed,
non-negative-score). A side-by-side positioning of the relevant
studies is in the replication package.

\paragraph{Contributions.}
\textbf{(C1) The \emph{strict} benign-pruning reading is rejected
at the within-tag value margin.} A single pre-specified shape
test (one model over the mutually exclusive net-score bins,
triple difference interacted with score level) returns a negative
average displacement across the value distribution ($-0.167$,
$p<10^{-4}$) attenuating with score (slope $p=2\times10^{-4}$);
the displacement is roughly constant across the low-to-mid range
($-0.21$, $-0.20$, $-0.21$ at bins $0,1,2$) and fades only at the
rare elite tail, so a channel confined to score-zero posts is
rejected---by a declared statistic, not a search over thresholds.
\textbf{(C2) A within-tag, within-question-type triple difference
identifies the LLM-substitutability slice.} At the
(tag $\times$ week $\times$ question-type) cell with
tag-by-question-type and week fixed effects, the baseline is
$-0.108$ ($p=0.021$), strengthening to $-0.140$ ($p=0.005$) once
the declared boundary category advanced-architecture is excluded;
a specification curve, placebos, a positive Copilot placebo,
PPML, two-way clustering, a wild bootstrap, and a HonestDiD
breakdown at $\bar M=0.75$ support the conditional interpretation.
\textbf{(C3) A value-calibrated early-warning system whose
detection signal is chosen by out-of-sample evidence.} On the
full 100-tag panel the activity-decline metric a platform already
tracks predicts realised 2025 high-value erosion at AUC $0.98$
versus $0.57$ for a value-weighted detector, so the system
\emph{detects} on activity while the value-margin estimates
\emph{calibrate} the alert into high-value posts at risk
(allocating by raw activity leaves $3$--$5\%$ of the protectable
flow unprotected); a persistent composition flag ($\rho=0.90$
into 2025) classifies intervention \emph{type}, drawing on a
five-transition funnel ($\sim$28{,}000-post shortfall, read as an
accounting device).

\paragraph{What this paper does not claim.}
We do not document a ``shock'': we estimate a marginal shortfall
of $\sim$70{,}000 questions over 109 weeks ($\sim$4\% of
post-period weekly volume). We do not document an effect on
innovation \emph{outcomes}, contributor entry, or productivity:
the Q\&A margin is an \emph{input} to open innovation. We do not
derive a welfare conclusion, since the private productivity gains
of LLM adoption
\citep{brynjolfsson2025generative,peng2023impact,cui2024effects}
are not estimated. The full robustness battery, mechanism,
validation, and positioning detail are in the public replication
package; \S\ref{sec:theory}--\S\ref{sec:conclusion} develop the
framework, data, identification, results, mechanisms, validation,
the early-warning system, limitations, and conclusion in turn.
